\documentclass[../../main.tex]{subfiles}
\begin{document}
% SOURCE p=238; printed=223; object=ch07-riesz-complex; number=7.17; region=upper
\begin{theorem}[Riesz 表示定理]\label{theorem:ch07-riesz-complex}
设 $X$ 为 LCH 空间。对 $\mu\in M(X)$ 及 $f\in C_0(X)$，令 $I_\mu(f)=\int f\,\mathrm d\mu$。则映射 $\mu\mapsto I_\mu$ 是从 $M(X)$ 到 $C_0(X)^*$ 的等距同构。
\end{theorem}
% SOURCE p=238; printed=223; object=ch07-riesz-complex-proof; region=upper-middle
\begin{proof}
我们已经证明每个 $I\in C_0(X)^*$ 都有 $I_\mu$ 的形式。另一方面，若 $\mu\in M(X)$，由\subcref{proposition:ch03-complex-variation-properties}{item:ch03-complex-integral-variation}，
\begin{align*}
\left|\int f\,\mathrm d\mu\right|\leqslant\int|f|\,\mathrm d|\mu|\leqslant\|f\|_u\|\mu\|,
\end{align*}
故 $I_\mu\in C_0(X)^*$ 且 $\|I_\mu\|\leqslant\|\mu\|$。又令 $h=\mathrm d\mu/\mathrm d|\mu|$，则由\subcref{proposition:ch03-complex-variation-properties}{item:ch03-complex-polar-density}，$|h|=1$ 几乎处处。根据 Lusin 定理，对任意 $\varepsilon>0$，存在 $f\in C_c(X)$，使 $\|f\|_u\leqslant1$，并且除去某个满足 $|\mu|(E)<\varepsilon/2$ 的集合 $E$ 外，$f=\overline h$。于是
\begin{align*}
\|\mu\|&=\int|h|^2\,\mathrm d|\mu|=\int\overline h\,\mathrm d\mu\\
&\leqslant\left|\int f\,\mathrm d\mu\right|+\left|\int(f-\overline h)\,\mathrm d\mu\right|\\
&\leqslant\left|\int f\,\mathrm d\mu\right|+2|\mu|(E)
<\left|\int f\,\mathrm d\mu\right|+\varepsilon\leqslant\|I_\mu\|+\varepsilon.
\end{align*}
所以 $\|\mu\|\leqslant\|I_\mu\|$，证明完成。
\end{proof}
% SUPPLEMENT: id=v14-ch07-riesz-complex-details; parent=ch07-riesz-complex; kind=proof-detail+motivation+hypotheses; origin=independent-derivation
\begin{remark}
\textbf{等距性、唯一性与条件（重构新增）。}动机是把任意连续线性“读数”解释成积分，并使测度的大小与读数的算子范数完全一致。正测度情形的质量最大化，在复测度情形变成抵消相位 $h$：理想测试函数是 $\overline h$，它可能不连续，故用\cref{theorem:ch07-lusin}在有限测度 $|\mu|$ 下替代。可取 $h$ 的 Borel 代表，并在零测集上改成 $1$，因而处处 $|h|=1$；这种修改不改变积分。有限总变差保证 Lusin 的有限支撑测度条件，也使误差至多 $2|\mu|(E)$。

原书此处写 $\|f\|_u=1$。证明只需上述 $\leqslant1$；若 $\mu\ne0$ 且先把误差取到小于 $\|\mu\|$，$f=\overline h$ 的集合有正质量，便确有 $\|f\|_u=1$。若 $\mu=0$，等距等式直接成立，无须寻找范数为 $1$ 的函数；这也包括 $X=\varnothing$。

积分对测度线性，所以此映射线性；等距性使其核为零，给出表示的唯一性；前面的四个正测度分解给出满射性。因此确为等距同构。由于对偶空间完备，$M(X)$ 也成为 Banach 空间：测度 Cauchy 列的积分泛函收敛到一个对偶元素，由满射性表示回来，再由等距性得到测度范数收敛。LCH 假设为紧支撑连续分离函数和正 Riesz 定理提供基础；Radon 正则性保证连续测试足以识别测度，一般 Borel 测度不具有这种唯一性，见\cref{exercise:ch07-ordinal-measure}。
\end{remark}
% END SUPPLEMENT: id=v14-ch07-riesz-complex-details
% SOURCE p=238; printed=223; object=ch07-riesz-compact; number=7.18; region=middle
\begin{corollary}\label{corollary:ch07-riesz-compact}
若 $X$ 为紧 Hausdorff 空间，则 $C(X)^*$ 与 $M(X)$ 等距同构。
\end{corollary}
% SUPPLEMENT: id=v14-ch07-riesz-compact-details; parent=ch07-riesz-compact; kind=omitted-proof+motivation+hypotheses; origin=independent-derivation
\begin{proof}
\textbf{原书省略证明的补全（重构新增）。}紧 Hausdorff 空间是 LCH 空间。每个连续函数在紧空间上有界，且对每个 $\varepsilon>0$，闭集 $\{|f|\geqslant\varepsilon\}$ 紧，所以每个 $f\in C(X)$ 都属于 $C_0(X)$。反向包含来自定义，故两者连同一致范数完全相同，应用前定理即得。
\end{proof}
\begin{remark}
\textbf{动机与条件（重构新增）。}紧性消除了无穷远边界，因而 $C_0$ 中的消失条件自动满足。Hausdorff 条件与紧性共同保证前定理所需的局部紧 Hausdorff 结构。非紧空间上的 $C(X)$ 中可能有无界函数，不能原样使用一致范数及本结论。
\end{remark}
% END SUPPLEMENT: id=v14-ch07-riesz-compact-details
% SOURCE p=238; printed=223; object=ch07-lone-measure-embedding; region=middle-lower
设 $\mu$ 是 $X$ 上一个固定的正 Radon 测度。若 $f\in L^1(\mu)$，则由 $\mathrm d\nu_f=f\,\mathrm d\mu$ 定义的复测度容易看出是 Radon 测度（\cref{exercise:ch07-integrable-weight-radon}），且 $\|\nu_f\|=\int|f|\,\mathrm d\mu=\|f\|_1$。因此 $f\mapsto\nu_f$ 是 $L^1(\mu)$ 到 $M(X)$ 的等距嵌入，其像恰为满足 $\nu\ll\mu$ 的 $\nu\in M(X)$。（最后一个断言来自 Radon--Nikodym 定理；即使 $\mu$ 不为 $\sigma$ 有限，它仍适用，见\cref{section:ch07-notes}。）最重要的例子是 $\mu=m$ 为 $\mathbb R^n$ 上的 Lebesgue 测度；我们将把 $L^1(m)$ 与 $M(\mathbb R^n)$ 的一个子空间等同。
% SUPPLEMENT: id=v14-ch07-lone-measure-embedding-details; parent=ch07-lone-measure-embedding; kind=omitted-proof+motivation+hypotheses; origin=independent-derivation
\begin{remark}
\textbf{包括非 $\sigma$ 有限情形的证明（重构新增）。}极分解给出 $|\nu_f|=|f|\mu$；\cref{exercise:ch07-integrable-weight-radon}使这个有限正测度 Radon，再由\cref{proposition:ch07-variation-radon}知 $\nu_f\in M(X)$。积分的线性和上面的范数等式给出等距线性映射；若两个 $L^1$ 等价类得到同一测度，则其差的 $L^1$ 范数为零，故它们相同。每个像中的测度显然绝对连续于 $\mu$。

反向设 $\nu\in M(X)$ 且 $\nu\ll\mu$。有限 Radon 测度 $|\nu|$ 在 $X$ 上内正则，故可取紧集 $K_n$ 使 $|\nu|(X\setminus K_n)<1/n$。令 $A=\bigcup_nK_n$，则 $|\nu|(X\setminus A)=0$；而 $\mu(K_n)<\infty$，所以 $\mu$ 在 Borel 子空间 $A$ 上的限制为 $\sigma$ 有限。对该限制应用\cref{theorem:ch03-complex-lebesgue-radon-nikodym}，得 $f_A\in L^1(\mu|_A)$，使 $\nu|_A=f_A\mu|_A$。把 $f_A$ 在 $X\setminus A$ 上延成零，所得 $f$ 可测可积；对任意 Borel 集 $E$，$\nu(E)=\nu(E\cap A)=\int_Ef\,\mathrm d\mu$，完成证明。这避免了把一般测度空间的 $\sigma$ 有限条件无理由删除。

这里 Radon 性提供了承载 $\nu$ 的可数个紧集，局部有限性使该承载集对 $\mu$ 也 $\sigma$ 有限；$L^1$ 条件恰保证有限总变差。动机是把函数看成“有密度的测度”，使极限能进入比 $L^1$ 更大的空间，例如集中成点质量。
\end{remark}
% END SUPPLEMENT: id=v14-ch07-lone-measure-embedding-details
% SOURCE p=238; printed=223; object=ch07-vague-topology-definition; region=lower
\begin{definition}[模糊拓扑与模糊收敛]\label{definition:ch07-vague-topology}
在 $M(X)=C_0(X)^*$ 上，把弱星拓扑称为\textbf{模糊拓扑}（vague topology）。其中网 $\mu_\alpha\to\mu$ 当且仅当对每个 $f\in C_0(X)$，
\begin{align*}
\int f\,\mathrm d\mu_\alpha\longrightarrow\int f\,\mathrm d\mu.
\end{align*}
这时称 $\mu_\alpha$ \textbf{模糊收敛}到 $\mu$。
\end{definition}
这种拓扑在应用中相当重要。“vague” 一词在概率论中很常见，而且在英语中比 “weak*” 更便于构成副词。模糊拓扑有时也叫弱拓扑，但这个叫法与本书的术语冲突，因为 $C_0(X)$ 很少是自反的（见\cref{exercise:ch07-czero-nonreflexive}）。对 $L^p(\mu)$ 的弱收敛论证一般在 $p=1$ 时失效，因为 $L^1(\mu)$ 并不是 $L^\infty(\mu)$ 的对偶；不过，经常可以像上段那样把 $L^1(\mu)$ 看作 $M(X)$ 的子空间，再利用后者的模糊拓扑获得合适的替代结论。
% SUPPLEMENT: id=v14-ch07-vague-topology-details; parent=ch07-vague-topology-definition; kind=motivation+hypotheses+proof-detail; origin=independent-derivation
\begin{remark}
\textbf{测试函数的选择（重构新增）。}本书明确使用全部 $C_0(X)$ 函数作测试，术语应以这一定义为准；仅用 $C_c(X)$ 测试的拓扑不能在任意无界测度族上未经证明地与之等同。若有共同范数界 $\sup_\alpha\|\mu_\alpha\|\leqslant M$，则 $C_c$ 测试收敛确实足够：对 $f\in C_0$ 取 $g\in C_c$ 一致逼近，
\begin{align*}
\left|\int f\,\mathrm d(\mu_\alpha-\mu)\right|
\leqslant(M+\|\mu\|)\|f-g\|_u+
\left|\int g\,\mathrm d(\mu_\alpha-\mu)\right|.
\end{align*}
先使第一项任意小，再沿网取极限即可。定义使用网以适应不一定可数可度量的弱星拓扑；对序列则直接使用同一测试条件。弱拓扑还要测试整个 $M(X)^*$，通常有更多测试泛函，因此可能更强。上段关于 $L^1$ 的话是一般情形的警示，不排除有限维等特殊空间存在额外对偶同构。
\end{remark}
% END SUPPLEMENT: id=v14-ch07-vague-topology-details
% SOURCE p=238; printed=223; object=ch07-vague-distribution-introduction; region=lower
本节最后给出 $M(\mathbb R)$ 中模糊收敛的一个有用判据。
\subfile{section-03-distribution.tex}
\end{document}
